\clearpage
\section{Literal clauses of the paper's term translation}
\label{sec:literal-translation}
\subsection{Fixed target and correspondence requirement}
The terms of Proposition~12.2 remain exactly
\[
 \mathtt{let}\ x=\mathtt{ref}\ 1\ \mathtt{in}\ \lambda f.(f();!x)
 \qquad\text{and}\qquad
 \lambda f.(f();1).
\]
Their existing scoped syntax and desugaring definitions are unchanged.
The formalisation must not replace either term by a different program, even
an intended equivalent one, to make the theorem easier. The translation
must follow Figure~16, including its canonical sharing trees, unused-resource
weakening, environment order, and box interfaces. In particular,
$D_{k,1}$ is not simplified to an identity: it contains a weakening and a
contraction for each resource.

This section implements the variable, atom, constant, unary, and abstraction
graph constructors literally as drawn in Figure~16. Section~\ref{sec:complete-translation}
now completes the application and binary clauses and the full recursive
term translation. Translation laws for substitution and contexts and the
observational equivalence proof remain open.

\subsection{Ordered binary juxtaposition}
\termdef{ordered-tensor}
For $G:I\to O$ and $H:J\to K$, ordered binary juxtaposition is a graph
$G\otimes H:(I\mathbin{+\!+}J)\to(O\mathbin{+\!+}K)$, placing all left
ports before all right ports. It is constructed by the previously proved
finite parallel assembly, rather than by adding identity wires.

A Boolean chooses the left or right interface. The port map sends a left
index $i$ to $i$, and a right index $j$ to $|I|+j$. Its inverse tests whether
the combined index is below $|I|$ and, in the right case, subtracts $|I|$.
The inverse equations follow by that case distinction and the index bounds;
they define a boundary bijection. Looking up a wire label in the concatenated
list gives the original component label, by the same left/right distinction.

\lean{Graph/Tensor}{9}{11}{Tensor.interfaces}

\lean{Graph/Tensor}{13}{15}{Tensor.join}

\lean{Graph/Tensor}{17}{19}{Tensor.split}

\lean{Graph/Tensor}{21}{26}{Tensor.split_join}

\lean{Graph/Tensor}{28}{30}{Tensor.join_split}

\lean{Graph/Tensor}{32}{33}{Tensor.ports}

\lean{Graph/Tensor}{35}{41}{Tensor.join_label}

The family selects $G$ or $H$. Apply the port bijections to both their
input and output interfaces to obtain the parallel boundary specification.
The tensor is the parallel construction for this two-element family, with
an explicit Boolean enumeration. Its graph laws are inherited from that
construction. Every tensor edge is a component tag and an original edge,
with unchanged label, so a label predicate holds on the tensor exactly
when it holds on both factors.

\lean{Graph/Tensor}{43}{45}{Tensor.family}

\lean{Graph/Tensor}{47}{51}{Tensor.boundary}

\lean{Graph/Tensor}{54}{55}{Tensor.net}

\lean{Graph/Tensor}{58}{64}{Tensor.labels_iff}

\subsection{Interface equalities and environment distributors}

An interface cast changes only provably equal type indices, leaving the
graph unchanged. Its input and output equalities are propositional equalities
of ordered wire lists, not a permission to permute ports. Eliminating these
equalities proves that every global primitive-label property is preserved
and reflected. Such casts reconcile, for example, $k\cdot1$ with $k$ and
$1+k$ with $k+1$; they perform no term rewrite or graph simplification.

\lean{Copying/TranslationClauses}{13}{14}{TranslationClauses.castInterfaces}

\lean{Copying/TranslationClauses}{16}{18}{TranslationClauses.allLabels_cast}

Juxtapose $D^{\star}_{k,m}$ with $D^{\diamond}_{h,m}$. This gives input
interface $\star^{km}\otimes\diamond^{hm}$ and output interface
$\star^k\otimes\diamond^h$, preserving the variables-before-atoms order.
The single-use constructor sets $m=1$ and uses only an interface-equality
cast. It retains every weakening and contraction in these distributors.

\lean{Copying/TranslationClauses}{21}{23}{TranslationClauses.environmentDistributor}

\lean{Copying/TranslationClauses}{26}{27}{TranslationClauses.singleUse}

\subsection{The five implemented Figure 16 clauses}
\termdef{literal-translation-clauses}

For resource $i$ among $n$ resources of sort $s$, construct the ordered
juxtaposition
\[
 D^s_{i,0}\ \otimes\ D^s_{1,1}\ \otimes\ D^s_{n-i-1,0}.
\]
It has one input of sort $s$ and $n$ outputs. The index bound $i<n$ proves
$i+1+(n-i-1)=n$, so the output cast preserves the exact order. Thus the
selected resource is in position $i$, while all other resources retain
weakening edges. The variable clause uses this construction at term sort,
then juxtaposes the weakening distributor for all ambient atoms. This is
the first diagram of Figure~16, with the two variable-environment portions
represented by the prefix and suffix of the selected finite index.

\lean{Copying/TranslationClauses}{30}{41}{TranslationClauses.selectedResource}

\lean{Copying/TranslationClauses}{44}{47}{TranslationClauses.variableTerm}

The atom clause keeps the instance edge. Serially connect that edge to the
selected atom-resource graph, and juxtapose weakening for all ambient term
variables on the left. The instance produces the atom input of its
$D^{\diamond}_{1,1}$; that tree is not omitted. Outputs remain variables
followed by atoms. For a constant, juxtapose variable weakening, the constant
corolla, and atom weakening in exactly the order shown in Figure~16.
Interface arithmetic checks the sole term input and the full ambient output
interface, including environments whose resources are unused.

\lean{Copying/TranslationClauses}{50}{54}{TranslationClauses.atom}

\lean{Copying/TranslationClauses}{57}{60}{TranslationClauses.constant}

For a unary operation, serially assemble the unary-operation corolla,
the supplied translated subterm, and the term/atom single-use distributors.
This follows the bottom-left diagram of Figure~16 for every unary operation,
including reference creation and dereference.

For abstraction, the supplied body has one extra term output. Place that
body in a box with one bound variable, $k$ free variables, and $h$ free atoms.
The bound variable is the first body output, as in the existing scoped
syntax. The equality of $k+1$ and $1+k$ reconciles the numeric type indices;
it does not reorder that body interface. The box has a $T^1(\star)$ input
and exposes only the free environment. Serially assemble the lambda corolla,
this box, and the two single-use distributors. Only this clause increases
the box-depth allowance. These are graph constructors; they do not replace
lambda or sequencing in the terms being formalised.

\lean{Copying/TranslationClauses}{63}{65}{TranslationClauses.unary}

\lean{Copying/TranslationClauses}{68}{70}{TranslationClauses.lambdaBox}

\lean{Copying/TranslationClauses}{73}{75}{TranslationClauses.abstraction}

\subsection{Global alphabet and absence of focus and holes}
\termdef{translation-language-graph}
The translation language permits operation, contraction, weakening, and
instance primitives. It excludes store, focus, and hole primitives.
A language graph satisfies this restriction recursively inside every box.
Store edges arise later from execution of reference creation, not from
translating the syntax of a reference-creation expression.

The primitive predicate explicitly enumerates the allowed constructors;
the graph predicate lifts it through all boxes. For any primitive-label
predicate, tensor closure follows from the tensor label equivalence. A
primitive corolla satisfies the lifted predicate whenever its sole primitive
label does, by inspecting the alphabet at zero or positive depth.

\lean{Copying/TranslationClauses}{78}{80}{TranslationClauses.LanguagePrimitive}

\lean{Copying/TranslationClauses}{82}{82}{TranslationClauses.LanguageGraph}

\lean{Copying/TranslationClauses}{84}{87}{TranslationClauses.allLabels_tensor}

\lean{Copying/TranslationClauses}{89}{92}{TranslationClauses.allLabels_primitive}

Every environment distributor is a language graph because each component
contains only contraction and weakening edges. Interface casting gives the
single-use case. Repeated tensor closure proves the selected-resource case.
The variable case adds only atom weakening; the atom case adds term weakening
and the permitted instance primitive; the constant case adds only its
permitted operation primitive. These arguments establish the global alphabet
property for every ambient scope, not just for closed terms.

\lean{Copying/TranslationClauses}{94}{96}{TranslationClauses.environmentDistributor_language}

\lean{Copying/TranslationClauses}{98}{99}{TranslationClauses.singleUse_language}

\lean{Copying/TranslationClauses}{101}{107}{TranslationClauses.selectedResource_language}

\lean{Copying/TranslationClauses}{109}{111}{TranslationClauses.variable_language}

\lean{Copying/TranslationClauses}{113}{116}{TranslationClauses.atom_language}

\lean{Copying/TranslationClauses}{118}{124}{TranslationClauses.constant_language}

The unary clause preserves the language property by serial closure,
combining its operation, the supplied subterm property, and the distributor
property. For the lambda box, the recursive label predicate is precisely
the body's label predicate, after the harmless interface cast. Serial closure
then proves the abstraction clause from the lambda primitive, that box,
and the single-use distributors. These are conditional preservation laws
for constructor arguments, not unproved assumptions about the final programs.

\lean{Copying/TranslationClauses}{126}{129}{TranslationClauses.unary_language}

\lean{Copying/TranslationClauses}{131}{134}{TranslationClauses.lambdaBox_language}

\lean{Copying/TranslationClauses}{136}{139}{TranslationClauses.abstraction_language}

A language graph is globally focus-free and hole-free. To prove each
claim, first lift the corresponding implication on primitive labels through
the recursive alphabet by induction on depth. At depth zero and for a
primitive at positive depth, equality with a forbidden label contradicts
the language predicate. For a box, apply the induction hypothesis to each
body edge. Apply this implication to every edge of the outer graph.

\lean{Copying/TranslationClauses}{142}{155}{TranslationClauses.language_focusFree}

\lean{Copying/TranslationClauses}{157}{169}{TranslationClauses.language_holeFree}

\subsection{Checks that preserve the paper's graph structure}

The reference-one graph assembles the literal constant and unary clauses
for the unchanged subterm $\mathtt{ref}\ 1$. It is a language graph in every
ambient environment by the two preservation lemmas. This is not the whole
left side of Proposition~12.2: the enclosing let and applications are assembled by the complete translation
in Section~\ref{sec:complete-translation}.

\lean{Copying/TranslationExamples}{7}{8}{TranslationExamples.referenceOne}

\lean{Copying/TranslationExamples}{10}{11}{TranslationExamples.referenceOne_language}

Kernel-checked finite-enumeration calculations guard the unsimplified
shapes. In an environment of two variables and one atom, a constant has
four outer edges: one constant and three weakenings. Selecting the second
of three variables with two ambient atoms gives six edges, including the
selected weakening/contraction pair. Selecting the second of two atoms
with three ambient variables gives seven, including the instance edge.
The graph for $\mathtt{ref}\ 1$ with two variables and one atom has eleven
edges: four for the constant clause, one for reference creation, and six
for the single-use environment distributors. These are structural checks,
not allocation or termination tests.

\lean{Copying/TranslationExamples}{14}{15}{TranslationExamples.constant_environment_edges}

\lean{Copying/TranslationExamples}{18}{19}{TranslationExamples.variable_environment_edges}

\lean{Copying/TranslationExamples}{22}{23}{TranslationExamples.atom_environment_edges}

\lean{Copying/TranslationExamples}{26}{27}{TranslationExamples.referenceOne_environment_edges}

The identity-abstraction regression supplies a variable graph at bound
index zero as the boxed body. Its language property follows from variable
and abstraction preservation. Its outer graph has exactly two edges, lambda
and box, when the free environment is empty; the variable's sharing graph
remains inside the box. Finally, the single-use distributors for two variables
and one atom have six edges. This directly checks the feature that would
be lost by replacing them with identity wiring.

\lean{Copying/TranslationExamples}{30}{31}{TranslationExamples.identityAbstraction}

\lean{Copying/TranslationExamples}{33}{34}{TranslationExamples.identityAbstraction_language}

\lean{Copying/TranslationExamples}{37}{37}{TranslationExamples.identityAbstraction_outer_edges}

\lean{Copying/TranslationExamples}{40}{41}{TranslationExamples.singleUse_retains_edges}

The full build and all three axiom audits pass with only standard Lean
axioms. The original proposition terms and their desugaring files have not
been modified. The outstanding translation and machine obligations remain
explicit in Appendix~\ref{app:remaining-work}; Proposition~12.2 is not yet proved.
